\chapter{Máquinas térmicas, entropía y segunda ley}
\label{cap:segunda_ley}

El cambio de entropía $dS$ entre dos estados de equilibrio infinitesimalmente separados es
$$ dS = \frac{dQ_r}{T} \quad (9.12) $$
donde $dQ_r$ es la transferencia de energía por calor para el sistema para un proceso reversible que conecta los estados inicial y final.
El cambio de entropía de un sistema durante un proceso finito arbitrario entre un estado inicial y un estado final es
$$ \Delta S = \int_i^f \frac{dQ_r}{T} \quad (9.13) $$
El valor de $\Delta S$ para el sistema es el mismo para todas las trayectorias que conectan los estados inicial y final. El cambio de entropía para un sistema que experimenta cualquier proceso cíclico reversible es cero.

\section*{Piense, dialogue y comparta}
Consulte el Prefacio para obtener una explicación de los iconos utilizados en este conjunto de problemas.

\begin{enumerate}
    \subimport{piense_01/}{ejercicio.tex}
    \subimport{piense_02/}{ejercicio.tex}
    \subimport{piense_03/}{ejercicio.tex}
\end{enumerate}

\section*{Problemas}

\subsection*{SECCIÓN 9.1 Máquinas térmicas y la segunda ley de la termodinámica}
\begin{enumerate}
    \subimport{prob_9_1_01/}{ejercicio.tex}
    \subimport{prob_9_1_02/}{ejercicio.tex}
    \subimport{prob_9_1_03/}{ejercicio.tex}
\end{enumerate}

\subsection*{SECCIÓN 9.2 Bombas de calor y refrigeradores}
\begin{enumerate}
    \setcounter{enumi}{3}
    \subimport{prob_9_2_04/}{ejercicio.tex}
    \subimport{prob_9_2_05/}{ejercicio.tex}
    \subimport{prob_9_2_06/}{ejercicio.tex}
\end{enumerate}

\subsection*{SECCIÓN 9.4 La máquina de Carnot}
\begin{enumerate}
    \setcounter{enumi}{6}
    \subimport{prob_9_4_07/}{ejercicio.tex}
    \subimport{prob_9_4_08/}{ejercicio.tex}
    \subimport{prob_9_4_09/}{ejercicio.tex}
    \subimport{prob_9_4_10/}{ejercicio.tex}
    \subimport{prob_9_4_11/}{ejercicio.tex}
    \subimport{prob_9_4_12/}{ejercicio.tex}
    \subimport{prob_9_4_13/}{ejercicio.tex}
    \subimport{prob_9_4_14/}{ejercicio.tex}
    \subimport{prob_9_4_15/}{ejercicio.tex}
    \subimport{prob_9_4_16/}{ejercicio.tex}
    \subimport{prob_9_4_17/}{ejercicio.tex}
\end{enumerate}

\subsection*{SECCIÓN 9.5 Motores de gasolina y diésel}
\begin{enumerate}
    \setcounter{enumi}{17}
    \subimport{prob_9_5_18/}{ejercicio.tex}
    \subimport{prob_9_5_19/}{ejercicio.tex}
\end{enumerate}

\subsection*{SECCIÓN 9.6 Entropía}
\begin{enumerate}
    \setcounter{enumi}{19}
    \subimport{prob_9_6_20/}{ejercicio.tex}
    \subimport{prob_9_6_21/}{ejercicio.tex}
\end{enumerate}

\subsection*{SECCIÓN 9.7 Entropía en sistemas termodinámicos}
\begin{enumerate}
    \setcounter{enumi}{21}
    \subimport{prob_9_7_22/}{ejercicio.tex}
    \subimport{prob_9_7_23/}{ejercicio.tex}
    \subimport{prob_9_7_24/}{ejercicio.tex}
    \subimport{prob_9_7_25/}{ejercicio.tex}
    \subimport{prob_9_7_26/}{ejercicio.tex}
\end{enumerate}

\subsection*{SECCIÓN 9.8 Entropía y la segunda ley}
\begin{enumerate}
    \setcounter{enumi}{26}
    \subimport{prob_9_8_27/}{ejercicio.tex}
    \subimport{prob_9_8_28/}{ejercicio.tex}
    \subimport{prob_9_8_29/}{ejercicio.tex}
\end{enumerate}

\section*{PROBLEMAS ADICIONALES}
\begin{enumerate}
    \setcounter{enumi}{29}
    \subimport{prob_adicional_30/}{ejercicio.tex}
    \subimport{prob_adicional_31/}{ejercicio.tex}
    \subimport{prob_adicional_32/}{ejercicio.tex}
    \subimport{prob_adicional_33/}{ejercicio.tex}
    \subimport{prob_adicional_34/}{ejercicio.tex}
    \subimport{prob_adicional_35/}{ejercicio.tex}
    \subimport{prob_adicional_36/}{ejercicio.tex}
    \subimport{prob_adicional_37/}{ejercicio.tex}
    \subimport{prob_adicional_38/}{ejercicio.tex}
    \subimport{prob_adicional_39/}{ejercicio.tex}
    \subimport{prob_adicional_40/}{ejercicio.tex}
    \subimport{prob_adicional_41/}{ejercicio.tex}
    \subimport{prob_adicional_42/}{ejercicio.tex}
    \subimport{prob_adicional_43/}{ejercicio.tex}
\end{enumerate}
